\section{Introduction}\label{intro:section}

\subsection{What collisions preserve}

A collision redistributes momentum and energy while preserving certain totals.
These conserved quantities determine how many parameters are needed to describe
equilibrium, and hence how many unknowns appear in a macroscopic equation.
The passage from a collision model to a fluid description therefore begins
with a concrete question: what information, beyond the known conservation
laws, can survive every collision?

We consider four-wave collisions. Write $k_0,k_1$ for the two incoming momenta
and $k_2,k_3$ for the outgoing momenta; the four entries are called the
\emph{legs} of a collision. For a function $\varphi$ of momentum, set
$\varphi_i=\varphi(k_i)$ and define the full collision difference
\begin{equation}\label{intro:collision-difference}
 \Delta\varphi=\varphi_0+\varphi_1-\varphi_2-\varphi_3.
\end{equation}
A function for which this difference vanishes almost everywhere on the allowed
resonances is a \emph{collision invariant}. Constants are always invariant,
and conservation of energy makes the dispersion relation invariant as well.
The question is whether these elementary examples exhaust the possibilities.

We follow the consequences of this question in two models. In a one-dimensional
pinned chain, resonance identities first force measurable invariants to be
smooth and then restrict them to constants and energy. In a three-dimensional
truncated model, five invariants determine the local equilibrium family.
Matching the actual five moments then makes microscopic and macroscopic errors
add exactly in relative entropy. The same collision difference thus links
classification to transport: the reciprocal of an equilibrium lies in its
kernel, while its weighted square measures dissipation away from equilibrium.

This relation is part of the programme of mechanical limits in Hilbert's
sixth problem. For hard spheres, Deng, Hani and Ma connected the long-time
derivation of the Boltzmann equation to fluid limits. Starting from elastic
hard-sphere dynamics with random initial configurations on two- and
three-dimensional tori, they obtained compressible Euler and incompressible
Navier--Stokes--Fourier systems under the corresponding initial-data and joint
scaling assumptions \cite[Theorems 1--3]{DHM2025}. Weakly nonlinear waves provide
another kinetic route. Deng and Hani derived the spatially homogeneous wave
kinetic equation from weakly nonlinear Schr\"odinger dynamics on large tori
\cite{DH2023}, and extended the derivation to any prescribed time for which the
kinetic solution exists \cite{DH2024}. Here we start from specified collision
equations and study their resonance geometry and macroscopic limits.
The dispersion, momentum cutoff and initial data are stated for each model.
Deriving these inhomogeneous collision models from microscopic wave dynamics
requires a separate analysis.

\subsection{Regularity from resonance compatibility}

The dimensionless dispersion of the nearest-neighbour pinned chain is
\begin{equation}\label{intro:pinned-dispersion}
 \omega_d(x)=\sqrt{1-2d\cos x},\qquad
 x\in\T_{2\pi}:=\R/(2\pi\mathbb Z),\qquad 0<d<\tfrac12.
\end{equation}
Here $d$ is a fixed coupling parameter and $x$ is the angular coordinate
of lattice momentum. The allowed collisions satisfy
\begin{equation}\label{intro:resonance}
 \begin{aligned}
 x_0+x_1&=x_2+x_3\pmod{2\pi},\\
 \omega_d(x_0)+\omega_d(x_1)&=\omega_d(x_2)+\omega_d(x_3).
 \end{aligned}
\end{equation}
Momentum is conserved modulo the period, including Umklapp processes across
the boundary of a fundamental interval. The angular coordinate $x$ itself
is therefore not a globally defined periodic collision invariant.

In their 2006 study of heat transport in this chain, Aoki, Lukkarinen and Spohn
asked whether constants and $\omega_d$ exhaust the collision invariants
\cite[after (4.10)]{ALS2006}. Lukkarinen later discussed almost-everywhere
solutions of this functional equation \cite[Section 3.2]{Luk2016}.
The review by Germain, La and Menegaki \cite[Section 7.2]{GLM2026} still lists
the classification for general $0<d<1/2$ as open; its Theorem 4 proves that
the space of $C^1$ solutions is finite-dimensional.

Theorem \ref{pin:classification} gives the classification. For every fixed
$0<d<1/2$, any almost-everywhere finite measurable collision invariant
$\varphi$ satisfies
\begin{equation}\label{intro:classification}
 \varphi=A+B\omega_d\quad\text{almost everywhere},\qquad A,B\in\C.
\end{equation}
The null sets in the collision identity are measured by the coarea measure
on the regular resonance surface; those in the conclusion are measured by
Lebesgue measure on the circle. Both measures are defined at the beginning
of Section \ref{pin:section}. No continuity or integrability is assumed.

The first step explains why a measurable invariant cannot remain arbitrarily
rough. For a target point $x$, choose a smooth resonance chart
$\bigl(x,H(x,z),z,W(x,z)\bigr)$, where $z$ is a free parameter and
$W=x+H-z$. The collision identity relates $\varphi(x)$ to its values at
three other points. These points move nondegenerately with $z$, so the
identity both controls the size of the function and permits averaging.
Three changes of variables first compare level sets and yield local
boundedness. A smooth function $\chi$, supported inside the chart and
having integral one, then gives, for almost every $x$,
\begin{equation}\label{intro:resonant-average}
 \varphi(x)=\int\chi(z)
 \bigl[\varphi(z)+\varphi(W(x,z))-\varphi(H(x,z))\bigr]\dd z.
\end{equation}
Since $\partial_zH$ and $\partial_zW$ stay away from zero, the right-hand side
can be written as an integral of $\varphi$ against a smooth kernel.
Derivatives in $x$ fall on the kernel, not on the measurable input.
The right-hand side defines a local smooth representative, and a finite
cover gives a smooth representative on the whole circle. Only then do we
differentiate the resonance identity. Second- and third-order compatibility,
together with periodic matching, completes the classification.
Regularity is a consequence of the collision relations between different
points.

\subsection{A common vanishing factor}\label{intro:shared-factor}

After identifying the kernel, one must give the dissipation integral a
suitable domain. The energy constraint degenerates at critical resonances,
and its induced measure may have infinite mass. The full collision
difference, however, vanishes at the same configurations.

The relation is visible in local coordinates. After permuting the legs,
write their positions as $(z+u,z+v,z,z+u+v)$. For a $C^2$ function $\varphi$,
\begin{equation}\label{intro:common-factor}
 \begin{aligned}
 F:=\Delta\omega_d&=-uvG_{\omega_d},&
 \Delta\varphi&=-uvG_\varphi,\\
 G_\varphi(z,u,v)&=\int_0^1\int_0^1
       \varphi''(z+su+tv)\dd s\dd t.
 \end{aligned}
\end{equation}
For example, the local test function $\varphi(z)=z^2$ gives
$\Delta\varphi=-2uv$. If $u=0$ or $v=0$, the legs coincide in pairs and
the difference vanishes. This quadratic function illustrates a local
calculation; it is not a periodic invariant on the circle.

When regularizing the energy constraint, we therefore first form
$|\Delta\varphi|^2=u^2v^2|G_\varphi|^2$ and then pass to the singular limit.
The square becomes small precisely where the degenerate measure becomes
large. Section \ref{pin:section} uses this common factor to construct a
densely defined closed dissipation form. The invariant classification
and compact resonance averages then give a positive spectral gap off
its kernel (Theorems \ref{pin:gap-theorem} and \ref{pin:closed-theorem}).
The gap constant depends on the fixed pinning parameter and positive
resonance weight.

\subsection{Moment matching and the decomposition of entropy}
\label{intro:entropy-spine}

The second model has a fixed momentum cube $D=[-R,R]^3$, with
$R>0$, a space variable $X\in\T^3:=\T_{2\pi}^3$, and time $t\ge0$.
Its dispersion is $|k|^2$, so the transport velocity is $v(k)=2k$.
We retain only collisions whose four momenta all lie in $D$.
Writing $\mathcal C$ for the resulting four-wave collision operator,
we study
\begin{equation}\label{intro:wke}
 \mathcal T f_c=c\mathcal C(f_c),\qquad
 \mathcal T:=\partial_t+2k\cdot\nabla_X,\qquad c\longrightarrow\infty.
\end{equation}
Here $c>0$ is the collision strength and $f_c(t,X,k)$ is a positive
distribution. Collisions act only on momentum: all four legs share the
same time and position. The measure, constants and cutoff in
$\mathcal C$ are specified in Section \ref{geom:section}.

The five collision invariants form the column vector
\[
 \Psi(k)=(1,k^{(1)},k^{(2)},k^{(3)},|k|^2)^{\mathsf T},
\]
where $k^{(j)}$ denotes the $j$th component of momentum. For a parameter
$\theta\in\R^5$ with $\theta\cdot\Psi>0$ throughout $D$, define the
Rayleigh--Jeans distribution, abbreviated as RJ,
\begin{equation}\label{intro:rj}
 N_\theta(k)=\frac1{\theta\cdot\Psi(k)}.
\end{equation}
Its reciprocal is a collision invariant, so $\mathcal C(N_\theta)=0$.
Allowing $\theta$ to depend on $(t,X)$ gives a local equilibrium.

If the five moments of a positive distribution $f$ lie in the image of
the equilibrium moment map, there is a unique $N=N_\theta$ such that
$\int_D\Psi(f-N)\dd k=0$. We call $N$ the equilibrium matched to the
actual five moments. For a reference equilibrium $\bar N=N_{\bar\theta}$,
define the relative entropy
\begin{equation}\label{intro:relative-entropy}
 \mathscr H(f\mid N)=\int_D
       \left[\frac fN-1-\log\frac fN\right]\dd k.
\end{equation}
The positive upper and lower bounds used below make this integral finite.
Moment matching gives the exact identity
\begin{equation}\label{intro:entropy-split}
 \boxed{\mathscr H(f\mid\bar N)
       =\mathscr H(f\mid N)+\mathscr H(N\mid\bar N).}
\end{equation}
Indeed, the cross term in the expansion is
\[
 \int_D(f-N)(\bar N^{-1}-N^{-1})\dd k
 =(\bar\theta-\theta)\cdot\int_D\Psi(f-N)\dd k=0.
\]
The first term measures the microscopic deviation that leaves the conserved
moments unchanged; the second measures the error in the macroscopic
parameters. They add exactly because the difference of the reciprocal
equilibria remains in the span of the five collision invariants.
Proposition \ref{hyd:entropy-geometry} gives the corresponding variational
properties, followed by an example with matched moments but positive
microscopic entropy.

\subsection{Corner memory and the Euler limit}

Evaluating the five conserved fluxes at local equilibrium gives the
macroscopic system
\begin{equation}\label{intro:euler}
 \partial_t\int_D\Psi N_\theta\dd k
 +\sum_{j=1}^3\partial_{X_j}\int_D2k^{(j)}\Psi N_\theta\dd k=0.
\end{equation}
It is symmetric hyperbolic in the positive parameter domain.
Let $\theta_0(X)$ be any smooth initial parameter with
$\theta_0\cdot\Psi$ uniformly positive, and suppose the Euler solution
$\bar\theta$ with this initial datum stays smooth and positive on a
finite interval $[0,T]$. We compare the kinetic solution starting from
$f_c(0)=N_{\theta_0}$ with $\bar N=N_{\bar\theta}$.
Write $N_c=N_{\theta_c}$ for the equilibrium matched to $f_c$.

Let $H_X^sL_k^2$ denote the Sobolev space with up to $s$ weak derivatives
in $X$ and values in $L^2(D)$, and take $L_t^\infty$ on $[0,T]$.
Theorem \ref{hyd:main} proves that, for sufficiently large $c$, a positive
classical solution exists up to $T$, with positive upper and lower bounds
independent of $c$, and
\begin{align}
 \|f_c-\bar N\|_{L_t^\infty H_X^3L_k^2}
 &\le C_Tc^{-3/4},\label{intro:kinetic-rate}\\
 \|\theta_c-\bar\theta\|_{L_t^\infty H_X^2}
 &\le C_Tc^{-1}.\label{intro:macro-rate}
\end{align}
The constant $C_T$ depends on finitely many derivatives of the reference
solution, its positivity margin and the fixed model parameters.
The spatial amplitude of the initial parameter is fixed; it need not
shrink as $c$ grows.

This limit allows momentum corners to retain a memory of the initial data.
If $\xi$ is a corner of the cube, the collision integral vanishes there,
and the positive classical solution satisfies
\begin{equation}\label{intro:corner-memory}
 f_c(t,X,\xi)=N_{\theta_0(X-2\xi t)}(\xi).
\end{equation}
The right-hand side generally differs from $\bar N(t,X,\xi)$.
Collisions are also slow near a corner: at distance $r$ from the corner
set, their frequency is of order $r^2\log^2(2R/r)$.
The proof must therefore control the contribution of the whole slow
region to dissipation and flux. The corners themselves have zero volume;
a thin corner layer has small volume, bounded amplitude, and still
collides with the rest of the momentum domain. Weighted energy estimates
that retain these interactions, together with amplitude bounds along
characteristics, allow corner memory to coexist with
\eqref{intro:kinetic-rate}.

Bricmont and Kupiainen \cite{BK2008} studied near-equilibrium diffusive
behaviour of a phonon Boltzmann equation, using wave numbers in at least
two dimensions, a modified dispersion without the square root, and a
diffusive scaling. Our one-dimensional classification retains the
square-root dispersion \eqref{intro:pinned-dispersion}. The three-dimensional
macroscopic theorem uses fixed transport and strong collisions, with
\eqref{intro:euler} as its limit. The resonance geometry required by each
model is proved separately.

\subsection{From dissipation to first-order transport}

The Euler equations use the equilibrium flux. The first-order correction
comes from the deviation from equilibrium. For a spatial direction
$i=1,2,3$ and a conserved quantity $a=0,\ldots,4$, define the residual flux
\begin{equation}\label{intro:remaining-flux}
 Z_{c,i}^a=\int_D2k^{(i)}\Psi_a(k)(f_c-N_c)\dd k.
\end{equation}
Theorem \ref{cur:main} proves, on the entire prescribed time interval,
\begin{equation}\label{intro:transport-limit}
 cZ_{c,i}^a\longrightarrow
 \sum_{j=1}^3\sum_{b=0}^4
 \mathsf K_{ij}^{ab}(\bar N)\partial_{X_j}\bar\theta_b
 \quad\text{in }L^2((0,T)\times\T^3).
\end{equation}
Here $\mathsf K$ is the Onsager transport tensor. Section \ref{ons:section}
defines it through the linearized collision equation after removing the
five zero modes, and proves that its entries are finite and that it is
symmetric and nonnegative. Symmetry of the collision form gives the
reciprocity relation $\mathsf K_{ij}^{ab}=\mathsf K_{ji}^{ba}$;
the general physical background is described in \cite{Onsager1931}.

The degeneracy of this tensor has a direct geometric description.
Write $\theta=(\alpha,\beta,\gamma)$, where $\alpha,\gamma\in\R$ and
$\beta\in\R^3$ correspond to the constant, momentum and energy invariants.
The fifteen spatial derivatives form a gradient array $\xi=(g,B,h)$,
where $g=\nabla_X\alpha$, $B_{jl}=\partial_{X_j}\beta_l$ and
$h=\nabla_X\gamma$. On every compact positive parameter set,
\begin{equation}\label{intro:onsager-degeneracy}
 \xi^{\mathsf T}\mathsf K(N_\theta)\xi
 \asymp
 \left|\frac{B+B^{\mathsf T}}2-\frac{\tr B}{3}I\right|^2+|h|^2.
\end{equation}
The symbol $\asymp$ denotes two-sided comparison with uniform positive
constants, and matrices carry the Frobenius norm.
The five-dimensional space of traceless symmetric $3\times3$ matrices
and the three energy-gradient directions account for the eight dissipative
directions. Mass gradients, skew-symmetric matrices and scalar matrices
account for the seven zero directions. At a nonzero spatial Fourier
frequency, gradient compatibility leaves only the mass direction.
Coupling to Euler transport then gives quantitative decay for the
constant-parameter linear system.

Strong convergence follows from the microscopic entropy balance under
actual moment matching: the change in microscopic entropy, four-wave
dissipation and work of the residual flux satisfy an exact identity.
A weak limit determines a candidate resonance current; the balance
also determines its norm and rules out a dissipation defect.
In addition to \eqref{intro:transport-limit}, we prove
\begin{equation}\label{intro:physical-current}
 c\mathcal C(f_c)\longrightarrow\mathcal T\bar N
 \quad\text{in }L^2((0,T)\times\T^3\times D).
\end{equation}
The full signed resonance current converges strongly with respect to
the resonance measure as well. These limits give the first-order transport
law for the actual five-moment parameters, with a remainder
$o(c^{-1})$ in $L^2(0,T;H^{-1}(\T^3))$, where $H^{-1}$ is the dual of
$H^1$. The convergence is integrated over the entire time interval;
first-order strong traces at the initial time require an initial-layer
analysis.

\subsection{Organization and conventions}

Section \ref{pin:section} proves the one-dimensional invariant classification
and constructs the closed dissipation form. Section \ref{geom:section}
turns to the three-dimensional model, establishing the pair-sum variance
representation, corner-frequency estimates and weighted coercivity.
Section \ref{hyd:section} uses the entropy decomposition to organize
nonlinear errors and prove the Euler limit. Section \ref{ons:section}
determines the transport tensor and its degeneracy, and
Section \ref{cur:section} uses the entropy balance to identify the strong
resonance current and all five moment fluxes.

Inner products are linear in the first argument. We write $\T_{2\pi}$ for
the circle and $\T^3=\T_{2\pi}^3$ for the spatial torus.
The measure $\dd X$ is periodic Lebesgue measure on $[0,2\pi)^3$,
and $\dd k$ is ordinary Lebesgue measure on the momentum domain.
Constants denoted by $C$ may depend on fixed model parameters; their
uniformity is specified in each statement. For $s>0$, we write
$\log_+s=\log^+s=\max\{\log s,0\}$.
Notation is introduced
where it is used and collected in an index at the end.
